\documentclass[../../../include/open-logic-section]{subfiles}

\begin{document}

\olfileid{sth}{ordinals}{opps}
\olsection{جانشین تے حدّی ترتیبی عدد}

اگلے کجھ باباں وچ اسیں ترتیبی عدد بہت ورتاں گے۔ ایس لئی کجھ سادے تصور
متعارف کرانے کم آؤن گے۔

\begin{defn}
کسے وی ترتیبی عدد $\alpha$ دا \emph{جانشین}
$\ordsucc{\alpha} =\alpha \cup \{\alpha\}$ اے۔
اسیں $\alpha$ نوں \emph{جانشین} ترتیبی عدد آکھدے آں جے
$\ordsucc{\beta} = \alpha$ کسے ترتیبی عدد~$\beta$ لئی۔
اسیں $\alpha$ نوں \emph{حدّی} ترتیبی عدد اُدوں تے صرف اُدوں آکھدے آں
جدوں $\alpha$ نہ خالی ہووے، نہ جانشین ترتیبی عدد۔
\end{defn}
\noindent
ہِیٹھلا نتیجہ وکھاندا اے کہ \emph{جانشین} دا ایہو ٹھیک تصور اے:

\begin{prop}
کسے وی ترتیبی عدد $\alpha$ لئی:
\begin{enumerate}
	\item $\alpha \in \ordsucc{\alpha}$۔
	\item $\ordsucc{\alpha}$ ترتیبی عدد اے۔
	\item کوئی ترتیبی عدد $\beta$ نہیں جیہدے لئی
	$\alpha \in \beta \in \ordsucc{\alpha}$۔
\end{enumerate}
\end{prop}

\begin{proof}
ظاہر اے کہ $\alpha \in \alpha \cup \{\alpha\} = \ordsucc{\alpha}$۔
ایسے طرح $\ordsucc{\alpha}$ ترتیبی عدداں دا متعدی سیٹ اے، تے ایس لئی
\olref[basic]{corordtransitiveord} نال ترتیبی عدد اے۔
تے $\alpha \in \beta \in \ordsucc{\alpha}$ ہونا ناممکن اے، کیوں جے
فیر یا $\beta \in \alpha$ یا $\beta = \alpha$، جیہڑا
\olref[basic]{ordordered} دے اُلٹ اے۔
\end{proof}
\noindent
ایس نال لامتناہی درجیاں اُتے استقرا نال ثبوت دا اک ہور طریقہ وی ملدا اے:

\begin{thm}[لامتناہی درجیاں اُتے سادہ استقرا]\ollabel{simpletransrecursion}
منّ لو $\phi(x)$ اک فارمولا اے جیہدے لئی:
\begin{enumerate}
	\item $\phi(\emptyset)$، تے
	\item کسے وی ترتیبی عدد $\alpha$ لئی، جے $\phi(\alpha)$، تاں
	$\phi(\ordsucc{\alpha})$، تے
	\item جے $\alpha$ حدّی ترتیبی عدد اے تے
	$(\forall \beta \in \alpha)\phi(\beta)$، تاں $\phi(\alpha)$۔
\end{enumerate}
فیر $\forall \alpha \phi(\alpha)$۔
\end{thm}

\begin{proof}
اسیں عکسِ نقیض ثابت کردے آں۔ سو منّ لو کوئی ترتیبی عدد اے جیہدے لئی
$\lnot\phi$۔ منّ لو $\gamma$ اوہ سبھ توں چھوٹا ترتیبی عدد اے۔
فیر یا $\gamma = \emptyset$، یا $\gamma = \ordsucc{\alpha}$ کسے ایہو
جیہے $\alpha$ لئی جیہدے لئی $\phi(\alpha)$، یا $\gamma$ حدّی ترتیبی عدد
اے تے $(\forall \beta \in \gamma)\phi(\beta)$۔
\end{proof}
\noindent
اک آخری نشان اگے کم آوے گا:

\begin{defn}\ollabel{defsupstrict}
جے $X$ ترتیبی عدداں دا سیٹ اے، تاں
$\supstrict(X) = \bigcup_{\alpha \in X} \ordsucc{\alpha}$۔
\end{defn}
\noindent
ایتھے ’’lsub‘‘ دا مطلب ’’سبھ توں چھوٹی سخت اُپَرلی حد‘‘ اے۔\footnote{کجھ
کتاباں ایس لئی ’’$\text{sup}(X)$‘‘ ورتدیاں نیں۔ پر ہور کتاباں
’’$\text{sup}(X)$‘‘ سبھ توں چھوٹی \emph{غیر سخت} اُپَرلی حد لئی
ورتدیاں نیں، یعنی بس $\bigcup X$۔ جے $X$ دا سبھ توں وڈا عنصر $\alpha$
ہووے، تاں ایہ تصور وکھرے ہو جاندے نیں۔ سبھ توں چھوٹی \emph{سخت} اُپَرلی
حد $\ordsucc{\alpha}$ اے، جدکہ سبھ توں چھوٹی \emph{غیر سخت} اُپَرلی حد
بس $\alpha$ اے۔} ہِیٹھلا نتیجہ ایس گل نوں سمجھاندا اے:

\begin{prop}
جے $X$ ترتیبی عدداں دا سیٹ اے، تاں $\supstrict(X)$ اوہ سبھ توں چھوٹا
ترتیبی عدد اے جیہڑا $X$ دے ہر ترتیبی عدد توں وڈا اے۔
\end{prop}

\begin{proof}
منّ لو $Y = \Setabs{\ordsucc{\alpha}}{\alpha \in X}$، سو
$\supstrict(X) = \bigcup Y$۔ چونکہ ترتیبی عدد متعدی نیں تے ترتیبی عدد
دا ہر رکن ترتیبی عدد اے، $\supstrict(X)$ ترتیبی عدداں دا متعدی سیٹ اے۔
سو \olref[basic]{corordtransitiveord} نال ایہ ترتیبی عدد اے۔

جے $\alpha \in X$، تاں $\ordsucc{\alpha} \in Y$، سو
$\ordsucc{\alpha} \subseteq \bigcup Y = \supstrict(X)$، تے ایس لئی
$\alpha \in \supstrict(X)$۔ سو $\supstrict(X)$، $X$ دے ہر ترتیبی عدد
توں سختی نال وڈا اے۔

اُلٹے پاسے، جے $\alpha \in \supstrict(X)$، تاں
$\alpha \in \ordsucc{\beta} \in Y$ کسے $\beta \in X$ لئی، سو
$\alpha \leq \beta \in X$۔ سو $\supstrict(X)$، $X$ دی \emph{سبھ توں
چھوٹی} سخت اُپَرلی حد اے۔
\end{proof}

\end{document}
